\documentclass[11pt]{article}

\usepackage[margin=1in]{geometry}
\usepackage{amsmath,amssymb,amsthm,mathtools}
\usepackage{microtype}

\newcommand{\Z}{\mathbb Z}
\newcommand{\R}{\mathbb R}
\newcommand{\one}{\mathbf 1}
\newcommand{\supp}{\operatorname{supp}}
\newcommand{\dist}{\operatorname{dist}}
\newcommand{\ip}[2]{\left\langle #1,#2\right\rangle}
\newcommand{\norm}[1]{\left\|#1\right\|}

\title{A Note on Equation (364) in Dimock's\\
\emph{The Renormalization Group According to Balaban, I. Small Fields}}
\author{}
\date{\today}

\begin{document}
\maketitle

\section*{Conclusion}

The estimate printed as Eq. (364),
\begin{equation}
  |C_k(\Omega;y,y')|
  \le O(1)L^2 \exp\{-\delta_0 L^{-2} d(y,y')\},
  \tag{364}
\end{equation}
is a valid weaker bound, but the factor \(L^{-2}\) is not the factor produced by
the scaling argument given immediately before it.  Scaling Eq. (348) from
\(G_{k+1}(L^{-1}\Omega)\) to \(G^0_{k+1}(\Omega)\) gives instead
\begin{equation}
  |\ip{f}{G^0_{k+1}(\Omega)f'}|
  \le C L^2 \exp\{-c L^{-1}d(y,y')\}\norm{f}_2\norm{f'}_2,
  \label{eq:sharp-g0}
\end{equation}
and hence
\begin{equation}
  |C_k(\Omega;y,y')|
  \le C L^2 \exp\{-c L^{-1}d(y,y')\}.
  \label{eq:sharp-ck}
\end{equation}
Here \(c>0\) is independent of \(k,N,\Omega,y,y'\), and may be taken as a fixed
constant multiple of the \(\delta_0\) in Eq. (348).

Since \(L^{-1}\ge L^{-2}\) for \(L>1\), the sharp estimate
\eqref{eq:sharp-ck} immediately implies the printed estimate (364), after
possibly changing constants.  Thus Eq. (364) is harmless for the later argument.
However, if Eq. (364) is read as the estimate that directly follows from scaling,
the exponent should contain \(L^{-1}\), not \(L^{-2}\).  This is also confirmed
internally by Eq. (370), whose first displayed line uses
\(\exp\{-\delta_0 L^{-1}d(y,y')\}\) while citing Eq. (364).

A further consequence is that the decay rate in Eq. (374) need not be of order
\(L^{-2}\).  If one uses \eqref{eq:sharp-ck} sharply, then the final estimate for
\(G_k(\Omega)\) may be written with an \(L\)-independent exponential rate,
\[
  |(G_k(\Omega)f)(x)|
  \le C_L\exp\{-\Gamma d(x,\supp f)\}\norm{f}_{\infty},
  \qquad \Gamma=O(1).
\]
Here \(C_L\) may depend on \(L\), as is allowed by the convention in the paper.
If, however, one keeps the extra weakening by \(L^{-1}\) that appears in the
printed Eq. (373), the resulting final rate is only \(O(L^{-1})\).  Thus the
printed \(O(L^{-2})\) in Lemma 33 is not forced by the improved Eq. (364).

The two substantive losses of sharpness are:
\[
  \hbox{Eq. (364): } L^{-1}\longrightarrow L^{-2},
  \qquad
  \hbox{Eq. (373): } \Gamma=L\alpha_L\longrightarrow \Gamma=\alpha_L.
\]
The first loss is the replacement of the scaling-correct \(C_k\) decay rate by a
weaker one.  The second loss is the extra \(L^{-1}\) inserted when Lemma 32 is
used inside the multiscale identity (372).

\section*{Notation}

In Appendix D, Dimock considers
\begin{equation}
  G_k(\Omega)
  =
  \left[-\Delta+\bar\mu_k+a_k Q_k^TQ_k\right]_{\Omega}^{-1},
  \tag{340}
\end{equation}
on \(\Omega\subset (L^{-k}\Z)^3\).  Equation (348) states that if
\(\supp f\subset \Delta_y\), \(\supp f'\subset \Delta_{y'}\), then
\begin{equation}
  |\ip{f}{G_k(\Omega)f'}|
  \le C\exp\{-\delta_0d(y,y')\}\norm{f}_2\norm{f'}_2.
  \tag{348}
\end{equation}

The operator relevant for Eq. (363) is
\begin{equation}
  G^0_{k+1}(\Omega)
  =
  \left[
    -\Delta+\bar\mu_k+L^{-2}a_{k+1}Q_{k+1}^TQ_{k+1}
  \right]_{\Omega}^{-1}.
  \tag{362}
\end{equation}
This is the unscaled \(k+1\) Green function acting on the old lattice
\((L^{-k}\Z)^3\).  It becomes the ordinary \(G_{k+1}\) after the spatial change of
variables \(x=Lz\).

\section*{Scaling of \(G^0_{k+1}\)}

Let
\[
  \Omega' = L^{-1}\Omega \subset (L^{-(k+1)}\Z)^3
\]
and define the pullback
\[
  (Rf)(z)=f(Lz), \qquad z\in \Omega'.
\]
With the usual weighted lattice sums, the measure scales by \(dx=L^3dz\), so
\begin{equation}
  \norm{Rf}_{L^2(\Omega')}=L^{-3/2}\norm{f}_{L^2(\Omega)}.
  \label{eq:norm-scale}
\end{equation}

The lattice Laplacian scales as
\begin{equation}
  R(-\Delta_k)R^{-1}=L^{-2}(-\Delta_{k+1}),
  \label{eq:lap-scale}
\end{equation}
and the mass satisfies
\begin{equation}
  \bar\mu_k=L^{-2}\bar\mu_{k+1}.
  \label{eq:mass-scale}
\end{equation}
Moreover the averaging term in \(G^0_{k+1}\) is precisely the old-lattice
representation of the \(k+1\) averaging term after scaling:
\begin{equation}
  R\left(L^{-2}Q_{k+1}^TQ_{k+1}\right)R^{-1}
  =
  L^{-2}Q_{k+1}^TQ_{k+1}.
  \label{eq:q-scale}
\end{equation}
The same symbol \(Q_{k+1}\) appears on both sides, but the left side acts on
\(\Omega\subset (L^{-k}\Z)^3\), whereas the right side acts on
\(\Omega'\subset (L^{-(k+1)}\Z)^3\).

Combining \eqref{eq:lap-scale}, \eqref{eq:mass-scale}, and \eqref{eq:q-scale},
\[
  R\left(G^0_{k+1}(\Omega)\right)^{-1}R^{-1}
  =
  L^{-2}\left(G_{k+1}(\Omega')\right)^{-1}.
\]
Therefore
\begin{equation}
  RG^0_{k+1}(\Omega)R^{-1}
  =
  L^2G_{k+1}(\Omega').
  \label{eq:operator-scale}
\end{equation}
Equivalently,
\[
  G^0_{k+1}(\Omega)f
  =
  L^2R^{-1}G_{k+1}(\Omega')Rf.
\]

Now compute the \(L^2\) pairing.  For any \(u\) on \(\Omega\) and \(v\) on
\(\Omega'\),
\[
  \ip{u}{R^{-1}v}_{\Omega}
  =
  L^3\ip{Ru}{v}_{\Omega'}.
\]
Using \eqref{eq:operator-scale},
\begin{align}
  \ip{f}{G^0_{k+1}(\Omega)f'}_{\Omega}
  &=
  L^2\ip{f}{R^{-1}G_{k+1}(\Omega')Rf'}_{\Omega} \notag\\
  &=
  L^5\ip{Rf}{G_{k+1}(\Omega')Rf'}_{\Omega'}.
  \label{eq:pairing-scale}
\end{align}

Apply Eq. (348) to \(G_{k+1}(\Omega')\).  The supports of \(Rf\) and \(Rf'\) are
the scaled sets \(L^{-1}\Delta_y\) and \(L^{-1}\Delta_{y'}\).  Their separation is
\[
  \dist(L^{-1}\Delta_y,L^{-1}\Delta_{y'})
  =
  L^{-1}d(y,y')+O(1),
\]
where the \(O(1)\) ambiguity only changes the prefactor.  Hence
\[
  |\ip{Rf}{G_{k+1}(\Omega')Rf'}_{\Omega'}|
  \le
  C\exp\{-cL^{-1}d(y,y')\}\norm{Rf}_2\norm{Rf'}_2.
\]
Using \eqref{eq:norm-scale} in \eqref{eq:pairing-scale} gives
\begin{align*}
  |\ip{f}{G^0_{k+1}(\Omega)f'}_{\Omega}|
  &\le
  L^5 C e^{-cL^{-1}d(y,y')}
    (L^{-3/2}\norm{f}_2)(L^{-3/2}\norm{f'}_2)\\
  &=
  C L^2 e^{-cL^{-1}d(y,y')}\norm{f}_2\norm{f'}_2.
\end{align*}
This proves \eqref{eq:sharp-g0}.  Thus the exponent obtained from scaling is
\(L^{-1}d(y,y')\), not \(L^{-2}d(y,y')\).

\section*{From \(G^0_{k+1}\) to \(C_k\)}

Equation (361) is
\begin{equation}
  C_k(\Omega)
  =
  A_k+a_k^2A_kQ_kG^0_{k+1}(\Omega)Q_k^TA_k,
  \tag{361}
\end{equation}
where
\[
  A_k=\left[a_k+aL^{-2}Q^TQ\right]_{\Omega}^{-1}.
\]
Since \(Q^TQ\) is a projection,
\[
  A_k
  =
  a_k^{-1}(I-Q^TQ)+(a_k+aL^{-2})^{-1}Q^TQ.
\]
Thus \(A_k\) is local, with uniformly bounded operator norms and bounded local
range.  The operators \(Q_k\) and \(Q_k^T\) are local averaging and extension
operators.  Consequently, \(A_kQ_k\) and \(Q_k^TA_k\) only smear a delta function
over a bounded neighborhood of the corresponding unit block, and their \(L^2\)
norms are uniformly bounded.

For the nonlocal term in (361), apply \eqref{eq:sharp-g0} to
\[
  f=Q_k^TA_k\delta_y,\qquad f'=Q_k^TA_k\delta_{y'}.
\]
The supports of these functions are within bounded distance of \(\Delta_y\) and
\(\Delta_{y'}\), respectively, and their \(L^2\) norms are \(O(1)\).  Therefore
\[
  |(A_kQ_kG^0_{k+1}(\Omega)Q_k^TA_k)(y,y')|
  \le
  C L^2 e^{-cL^{-1}d(y,y')}.
\]
The first term \(A_k(y,y')\) is local.  If it is nonzero, then \(d(y,y')=O(1)\),
and it is also bounded by the same right side after increasing \(C\).  Hence
\[
  |C_k(\Omega;y,y')|
  \le
  C L^2 e^{-cL^{-1}d(y,y')},
\]
which is \eqref{eq:sharp-ck}.

\section*{Consequence for Lemma 32}

Now insert the improved estimate \eqref{eq:sharp-ck} into the proof of Lemma 32.
Recall
\[
  C'_k(\Omega)=H_k(\Omega)C_k(\Omega)H_k(\Omega)^T.
\]
Lemma 31 gives, with \(\delta_1=O(1)\),
\[
  |H_k(\Omega;x,y)|\le C e^{-\delta_1d(x,y)},
  \qquad
  |H_k(\Omega;y',z)|\le C e^{-\delta_1d(y',z)}.
\]
Thus, for \(A=\supp f\),
\begin{align}
  |(C'_k(\Omega)f)(x)|
  &\le
  C_L
  \sum_{y,y'}
  e^{-\delta_1d(x,y)}
  e^{-cL^{-1}d(y,y')}
  e^{-\delta_1d(y',A)}
  \norm{f}_{\infty}.
  \label{eq:cprime-convolution}
\end{align}
The prefactor \(C_L\) absorbs the \(L^2\) in \eqref{eq:sharp-ck} and the
summability constants.  Since
\[
  d(x,A)\le d(x,y)+d(y,y')+d(y',A),
\]
the convolution in \eqref{eq:cprime-convolution} decays at the slowest of the
three displayed rates.  Therefore
\begin{equation}
  |(C'_k(\Omega)f)(x)|
  \le C_L e^{-\alpha_Ld(x,\supp f)}\norm{f}_{\infty},
  \qquad
  \alpha_L=c_0L^{-1},
  \label{eq:cprime-improved}
\end{equation}
with \(c_0=O(1)\) independent of \(L\).  The same reasoning gives the derivative
and Holder-derivative estimates, since the corresponding bounds on \(H_k\) have
the same \(O(1)\) decay rate \(\delta_1\).

Thus the improved Eq. (364) improves Lemma 32 from the paper's
\(\gamma_0=O(L^{-2})\) to an intermediate one-step rate \(O(L^{-1})\).  It does
not make Lemma 32 itself \(O(1)\), because \(C_k\) still has only
\(e^{-cL^{-1}d}\) decay.

\section*{Consequence for Eq. (374)}

The key point is that Eq. (372) contains an additional scaling by
\(L^{k-j}\).  Write \(n=k-j\), so \(1\le n\le k\), and put
\[
  D=d(x,\supp f).
\]
The identity (372) is
\begin{equation}
  (G_k(\Omega)f)(x)
  =
  \sum_{j=0}^{k-1}
  L^{-2(k-j)}
  (C'_j(L^{k-j}\Omega)f_{L^{k-j}})(L^{k-j}x).
  \tag{372}
\end{equation}
Equivalently,
\[
  (G_k(\Omega)f)(x)
  =
  \sum_{n=1}^{k}
  L^{-2n}
  (C'_{k-n}(L^n\Omega)f_{L^n})(L^nx).
\]
The support rescales as
\[
  d(L^nx,\supp f_{L^n})=L^n d(x,\supp f)=L^nD.
\]
Applying \eqref{eq:cprime-improved} directly gives
\begin{align}
  |(G_k(\Omega)f)(x)|
  &\le
  C_L
  \sum_{n=1}^{k}
  L^{-2n}e^{-\alpha_LL^nD}\norm{f}_{\infty}.
  \label{eq:g-sum-improved}
\end{align}
Since \(\alpha_L=c_0L^{-1}\),
\[
  \alpha_LL^n=c_0L^{n-1}\ge c_0
  \qquad (n\ge1).
\]
Hence each term in \eqref{eq:g-sum-improved} satisfies
\[
  e^{-\alpha_LL^nD}\le e^{-c_0D},
\]
and
\[
  \sum_{n=1}^{k}L^{-2n}\le \sum_{n=1}^{\infty}L^{-2n}
  \le C.
\]
Therefore the final Green-function estimate can be written as
\begin{equation}
  |(G_k(\Omega)f)(x)|
  \le
  C_L e^{-c_0d(x,\supp f)}\norm{f}_{\infty}.
  \label{eq:g-final-o1}
\end{equation}
The decay rate \(c_0\) is \(O(1)\), independent of \(L\).  The prefactor may
depend on \(L\), in agreement with the paper's convention that constants denoted
\(C\) may depend on \(L\).

This conclusion uses \eqref{eq:cprime-improved} sharply.  The printed Eq. (373)
contains an additional weakening:
\[
  e^{-\gamma_0 d(L^nx,\supp f_{L^n})}
  \quad\hbox{is replaced by}\quad
  e^{-\gamma_0 L^{-1} d(L^nx,\supp f_{L^n})}.
\]
If one makes the analogous weakening after improving Lemma 32, then
\[
  e^{-\alpha_L L^{-1}d(L^nx,\supp f_{L^n})}
  =
  e^{-c_0L^{n-2}D},
\]
and the first term \(n=1\) gives only \(e^{-c_0L^{-1}D}\).  In that weakened
reading the final rate is \(O(L^{-1})\), not \(O(1)\).

\section*{Where Sharpness Is Lost}

It is useful to separate the rates in the three objects
\[
  C_k,\qquad C'_k=H_kC_kH_k^T,\qquad G_k.
\]
Let
\[
  \beta_L=\hbox{decay rate in }C_k,\qquad
  \alpha_L=\hbox{decay rate in }C'_k,\qquad
  \Gamma_L=\hbox{final decay rate in }G_k.
\]
The sharp propagation of rates is
\[
  \beta_L
  \quad\Longrightarrow\quad
  \alpha_L\simeq \min\{\delta_1,\beta_L\}
  \quad\Longrightarrow\quad
  \Gamma_L\simeq L\alpha_L.
\]
The first implication is just the convolution with the two \(H_k\) kernels.  Since
the \(H_k\) kernels decay at the \(O(1)\) rate \(\delta_1\), they cannot improve a
slower \(C_k\) rate; they preserve it up to harmless constants.  The second
implication comes from the factor
\[
  d(L^nx,\supp f_{L^n})=L^n d(x,\supp f)
\]
in Eq. (372), and the worst case \(n=1\) gives the gain \(L\).

Thus, if the improved Eq. (364) is used, then
\[
  \beta_L\simeq L^{-1},\qquad
  \alpha_L\simeq L^{-1},\qquad
  \Gamma_L\simeq L\alpha_L\simeq O(1).
\]

The first non-sharp step is the printed scaling estimate (363), and hence the
printed Eq. (364).  Scaling Eq. (348) from \(G_{k+1}(L^{-1}\Omega)\) gives
\[
  |\ip{f}{G^0_{k+1}(\Omega)f'}|
  \le
  C L^2 e^{-cL^{-1}d(y,y')}\norm{f}_2\norm{f'}_2,
\]
because distances are divided by \(L\) under \(\Omega\mapsto L^{-1}\Omega\).
The printed \(e^{-\delta_0L^{-2}d(y,y')}\) is therefore a weakening by one power
of \(L\).  Passing from \(G^0_{k+1}\) to \(C_k\) through
\[
  C_k=A_k+a_k^2A_kQ_kG^0_{k+1}Q_k^TA_k
\]
does not force any further loss, because \(A_k,Q_k,Q_k^T\) are local.  Their
support enlargement is \(O(L)\), and
\[
  e^{-cL^{-1}(d-O(L))}\le C e^{-cL^{-1}d}.
\]

The second non-sharp step is Eq. (373).  If Lemma 32 gives
\[
  |(C'_jF)(X)|\le C_Le^{-\alpha_Ld(X,\supp F)}\norm{F}_\infty,
\]
then a sharp substitution into Eq. (372) gives
\[
  |(C'_{k-n}(L^n\Omega)f_{L^n})(L^nx)|
  \le
  C_Le^{-\alpha_LL^nd(x,\supp f)}\norm{f}_\infty.
\]
The printed Eq. (373) instead replaces this by
\[
  C_Le^{-\alpha_LL^{-1}d(L^nx,\supp f_{L^n})}\norm{f}_\infty
  =
  C_Le^{-\alpha_LL^{n-1}d(x,\supp f)}\norm{f}_\infty.
\]
This loses exactly one power of \(L\) in the final rate.  It changes
\[
  \Gamma_L=L\alpha_L
  \quad\hbox{into}\quad
  \Gamma_L=\alpha_L.
\]
Everything after this point is sharp relative to the weakened input; the
geometric sum in Eq. (374) does not introduce a new decay loss.

The complete rate ledger is therefore:
\[
\begin{array}{c|c}
\hbox{input/procedure} & \hbox{final rate in Eq. (374)}\\
\hline
\hbox{printed Eq. (364), used sharply in Eq. (372)} & O(L^{-1})\\
\hbox{printed Eq. (364), plus printed Eq. (373) weakening} & O(L^{-2})\\
\hbox{improved Eq. (364), but printed Eq. (373) weakening} & O(L^{-1})\\
\hbox{improved Eq. (364), used sharply in Eq. (372)} & O(1)
\end{array}
\]

\section*{Why the Printed \(L^{-2}\) Is Harmless}

The sharper estimate implies the printed estimate because, for \(L>1\) and
\(d\ge0\),
\[
  e^{-cL^{-1}d}\le e^{-cL^{-2}d}.
\]
Thus
\[
  C L^2e^{-cL^{-1}d(y,y')}
  \le
  C L^2e^{-cL^{-2}d(y,y')}.
\]
This explains why the later proof can safely use a decay rate of order \(L^{-2}\).
Indeed Lemma 32 explicitly sets
\[
  \gamma_0=\frac12\delta_0L^{-2}.
\]

However, Eq. (370) begins with
\[
  e^{-\delta_1d(x,y)}
  e^{-\delta_0L^{-1}d(y,y')}
  e^{-\delta_1d(y',\supp f)},
\]
while saying ``By (364) and (366).''  This line is incompatible with the printed
form of Eq. (364), but it is exactly compatible with the scaled estimate
\eqref{eq:sharp-ck}.  The most coherent reading is therefore:
\begin{itemize}
  \item the estimate actually obtained by scaling, and apparently used in the
  first line of Eq. (370), has \(L^{-1}\) in the exponent;
  \item the weaker \(L^{-2}\) bound is sufficient for the subsequent argument and
  is the rate finally retained in Lemma 32 and Lemma 33;
  \item if the improved estimate is propagated sharply through Eq. (372), the
  final decay rate in Eq. (374) can be taken independent of \(L\).
\end{itemize}

\end{document}
